\section{Sample Prompts}
\label{appendix:sample-prompts}

This appendix provides sample prompts you may use when interacting with an LLM.
You are not required to use these prompts verbatim.
However, you should keep the specification fixed and avoid importing external tax rules beyond what is stated in~\cref{appendix:tax-spec}.

\subsection{Prompt 1: Increasing Design Diversity}
\label{appendix:prompt-diversity}

\begin{tcolorbox}[colback=purple!10,colframe=purple!70!black,boxrule=0.5pt]
\textbf{Prompt:} Hi ChatGPT, I'm interested in developing diverse software implementations of a spec.
I am thinking about ways to prompt you that will increase the variation in the implementations.
I am wondering about programming language and library selection (vs.\ custom implementations of components).
Are there other things I should think about?
Please give concrete strategies and examples, and explain how each strategy might affect the diversity of failure modes.
\end{tcolorbox}

\subsection{Prompt 2: Starting from a System Design}
\label{appendix:prompt-system-design}

\begin{tcolorbox}[colback=purple!10,colframe=purple!70!black,boxrule=0.5pt]
\textbf{Prompt:} The full spec is available as a file in your context.
Can we start with a system design?
I want the design to have an NDJSON parser, an analysis engine, and an output writer.
I would also like to use a database to store intermediate computations, so that the tax computation pipeline is explicit and debuggable.
Please propose a modular architecture and data model first, before writing any code.
Then, after I confirm, we will implement the modules one at a time.
\end{tcolorbox}